\documentclass[10pt,a4paper]{article}
\usepackage[fleqn]{amsmath}
\usepackage{amssymb,amsthm,graphicx}
\usepackage{latexsym}
\usepackage[fleqn]{amsmath}
\usepackage{amssymb}
\usepackage{amsthm}
\usepackage{enumitem}
\newcommand{\HCF}{\mathrm{hcf}}
\newcommand{\ZZ}{\mathbb{Z}}
\newcommand{\rk}{\mathrm{rk}}
\DeclareMathOperator{\len}{length}
\DeclareMathOperator{\bl}{b}
\DeclareMathOperator{\cl}{cl}
\DeclareMathOperator{\di}{d}
\DeclareMathOperator{\Cl}{Cl}
\DeclareMathOperator{\acl}{acl}
\DeclareMathOperator{\tp}{tp}
\DeclareMathOperator{\Th}{Th}
\newcommand{\Cor}{\mathcal{C}_\mathrm{or}}
\newcommand{\Cun}{\mathcal{C}_\mathrm{un}}
\newcommand{\mons}{\mathbb{M}}
\newcommand{\strong}{\mathrm{strong}}

\newtheoremstyle{note}% hnamei
{7pt}% hSpace abovei
{10pt}% hSpace belowi
{\normalfont}% hBody fonti
{}% hIndent amounti1
{\bfseries}% hTheorem head fonti
{ }% hPunctuation after theorem headi
{.5em}% hSpace after theorem headi2
{}% hTheorem head spec (can be left empty, meaning ‘normal’)i
\theoremstyle{note}
\newtheorem{thm}{Theorem}[section]
%\newtheorem{theorem}{Theorem}
\newtheorem{lemma}[thm]{Lemma} 
\newtheorem{df}[thm]{Definition}
\newtheorem{rmk}[thm]{Remark} 
\newtheorem{prop}[thm]{Proposition}
\newtheorem{fact}[thm]{Fact}
\newtheorem{cor}[thm]{Corollary}
\newtheorem{conj}[thm]{Conjecture}
\newtheorem{ex}[thm]{Example}
\newtheorem{qu}[thm]{Question}
\newtheorem*{notation}{Notation} 
\title{Finitely generated unconfined Steiner triple systems are free}
\author{Silvia Barbina, Enrique Casanovas}
\date{arxiv:2412.06055v3; CC BY 4.0}
\usepackage{url}
\begin{document}
\maketitle
\noindent\textit{Editorial note.} The finite HF-ordering argument supplies nonnegative predimension by assigning each block to its greatest vertex, with at most one block per vertex. The cited companion paper is not granted as a hidden freeness theorem.\par\medskip
\noindent\textit{Editorial note.} The strict-order compactness display294 lacks a!=b in its totality clause. Only the finite peeling direction of this support proposition is needed here; the infinite compactness display is retained as source context, not used to infer the target.\par\medskip
\noindent\textit{Editorial scope.} One original proposition with its full author local proof chain; no new claim assigned to any support lemma.\par\medskip
\noindent\textit{Allowed standard prerequisites (editorial).}\begin{itemize}
\item Elementary universal algebra and well-founded recursion. Universal mapping property, generated subalgebras, uniqueness of a homomorphism on generators, and recursion/induction on well-orders.
\item Finite combinatorial accounting on terms and triples. Term length and rank, counting blocks by their greatest vertex, and monotone nonnegative integer sequences.
\end{itemize}\medskip
A \textbf{Steiner triple system} (STS) of order $\lambda$ is a pair $(M, \mathcal{B})$ where $M$ is a set with $\lambda$ elements, and 
 $\mathcal{B}$ is a collection of 3-element subsets of $M$ (the \textbf{blocks})
such that any two $x, y \in M$ are contained in exactly one block. If we require that any two $x,y \in M$ are contained in \textit{at most }one block, $M$ is said to be a  \textbf{partial Steiner triple system}.

From a model theoretic standpoint, Steiner triple systems can be viewed as relational structures with a single a ternary relation $R$ interpreted as the set of blocks (that is, three elements $x, y$ and $z$ are in a block if and only if $R(x,y,z)$). Then $R$ must satisfy
\begin{itemize}
\item for all permutations $\sigma \in S_3$, $ \forall x_1 x_2 x_3 \, \left(R(x_1, x_2, x_3) \rightarrow R(x_{\sigma(1)}, x_{\sigma(2)}, x_{\sigma(3)} \right)$
\item $\forall x y z \left( R(x,y,z) \rightarrow y\neq z \right)$
\item $\forall x y w z \left( R(x,y,w) \wedge R(x,y,z) \rightarrow w=z \right)$.
\end{itemize}
With this choice of language, a substructure of a Steiner triple system is a partial Steiner triple system.

An STS $(M, \mathcal{B})$ determines a \textit{Steiner quasigroup}, that is, a pair $(M, \cdot)$ where $\cdot$ is a binary operation such that
\begin{itemize}
\item $\forall x \, (x\cdot x =x)$
\item if $x \neq y$, then $x \cdot y = z$ where $\{x,y,z \} \in \mathcal{B}$.
\end{itemize}
The operation $\cdot$ defined in this way satisfies
 \begin{enumerate}
\item $\forall x \, (x \cdot y= y\cdot x)$
\item $\forall x \, (x\cdot x =x)$
\item $\forall x \, (x\cdot (x\cdot y)=y)$.
\end{enumerate}
Conversely, any structure $(M, \cdot)$ where $\cdot$ satisfies properties 1, 2 and 3 determines an STS $(M, \mathcal{B})$.

In universal algebraic parlance, Steiner quasigroups are algebras. Thus, if $M$ is a Steiner quasigroup, there are standard notions of a \textit{subalgebra} of $M$, and of the subalgebra \textit{generated} by a subset $A \subseteq M$. A subalgebra of a Steiner quasigroup is a substructure in the model theoretic sense, and in particular it is a Steiner quasigroup. If $A \subseteq M$, then $\langle A \rangle_M$ denotes the substructure of $M$ generated by $A$. The subscript $M$ is usually omitted when clear  from the context.

A Steiner quasigroup induces a partial STS structure on each of its subsets, so we often refer to a partial Steiner triple system when we need to consider a subset of a Steiner quasigroup that is not a substructure.

Given a class  $\mathcal{V}$ of algebras in a given signature, we say that a member $U$ of $\mathcal{V}$ generated by a subset $X \subseteq U$ has the  \textbf{universal mapping property} (UMP) over $X$ with respect to $\mathcal{V}$ if for every $V \in \mathcal{V}$, every map
 \[ \varphi : X \to V \]
extends to a homomorphism $ \hat{\varphi} : U \to V $. A standard argument proves that the homomorphism $\hat{\varphi}$ is unique. Such an algebra $U$ is said to be \textbf{freely generated by $X$}, and $X$ is a \textbf{set of free generators} of $U$. 

It is well known that when $\mathcal{V}$ is a variety, given any cardinal $\lambda \geq 1$ there are $U \in \mathcal{V}$ and $X \subseteq U$ such that $U$ is freely generated by $X $ and $|X|= \lambda$.

Since the class $\mathcal{K}$ of Steiner quasigroups is a variety in the universal algebraic sense, for each cardinal $\lambda \geq 1$, there is a Steiner quasigroup that is freely generated by a set of cardinality $\lambda$. Then a \textbf{free Steiner quasigroup} is a Steiner quasigroup that is freely generated by some subset $X$, and so it has UMP over $X$ with respect to the class~$\mathcal{K}$. We assume $\lambda \geq 3$ to exclude the trivial cases of Steiner quasigroups with one and three elements.

	\begin{df} Let $(M,\cdot)$ be a Steiner quasigroup. A subset $A\subseteq M$ is  \textbf{independent} if the generated substructure $\langle A\rangle_M$ is freely generated from $A$, that is, if $\langle A\rangle_M$ has UMP over $A$ for the class of Steiner quasigroups. 
		
		We say that $A\subseteq M$ is a \textbf{free base} from $M$  if it is independent and it generates $M$. 
	\end{df}
	When we say $\bar{a} = \langle a_i : i \in I \rangle$ is independent, we understand $a_i\neq a_j$ if  $i,j\in I$ and $i\neq j$.
	
	If $M$ is a Steiner quasigroup generated by the subset $A$ and $\bar{a}$ is an enumeration of $A$, then for every $b \in M$ there is a term $t(\bar{x})$ such that $b=t(\bar{a})$.

 Let $M$ be a Steiner quasigroup generated by $A\subseteq M$.  Then  $M=\bigcup_{n<\omega} S_n$ for a chain of sets $S_n$ such that  $S_0 =A$ and $S_{n+1}= \{a\cdot b\mid a,b\in S_n\}$. Note that $S_n\subseteq S_{n+1}$  since  $a\cdot a = a$.
 We will write $S_{<n}$ for $\bigcup_{i<n} S_i$. 
	 The sets $S_n$  are called the \textbf{levels} of $M$ over $A$. An element $a\in M$ has   \textbf{level~$n$ over $A$}  if $a\in S_n\smallsetminus S_{<n}$.

	\begin{df} A partial  STS $A$ is \textbf{unconfined}, or \textbf{contains no confined configuration}, if for all finite non-empty $A^\prime \subseteq A$ there is $a \in A^\prime$ that belongs to at most one block in $A^\prime$, that is, $a$ is the product of at most two elements of $A^\prime$.
		
	\end{df}

\begin{df} Let $A$ be a finite partial STS and let $\bl(A)$ be the set of blocks in $A$. Then the \textbf{rank}, or \textbf{predimension}, of $A$ is
	\[ \delta(A)= |A|-|\bl(A)| \, .\]
\end{df}

It is easy to see that the predimension of a finite Steiner quasigroup is negative except in the case of the Steiner quasigroup of order 7 (the Fano plane), which has predimension 0. On the other hand, the predimension of a finite unconfined partial Steiner triple system is always non-negative (see \cite{barcas3} for details).

 Let $Aa$ be a partial Steiner triple system, with $a\not\in A$. Then 
$$\delta(Aa)= \delta(A)+1-|\bl_a(A)| \, ,$$
 where $\bl_a(A)$ is the set of blocks containing $a$ and intersecting $A$. Hence, if there is at most one such block, $\delta(Aa)\geq \delta(A)$. It follows that if $A\subseteq B$ are finite partial Steiner triple systems and there is an HF-ordering of $B$ with initial segment $A$, then $\delta(B)\geq\delta(A)\geq 0$.

\begin{prop}\label{2024_nov4}
	Finitely generated unconfined Steiner triple systems are free Steiner quasigroups.
\end{prop}

	\begin{prop}\label{May12_2} Let $M$ be a Steiner quasigroup and let $A\subseteq M$. The following are equivalent:
		\begin{enumerate}
			\item $A$ is a free base of $M$.
			\item $M$ is generated by $A$ with  levels $(S_n\mid n<\omega)$ such that $A=S_0$ and
			\begin{enumerate}
				\item if $a,b\in S_n\smallsetminus S_{<n}$ and $a\neq b$, then $a\cdot b\not\in S_n$.
				\item if $a\in S_{n+1}\smallsetminus S_n$, there is a unique pair $\{b,c\}\subseteq S_n$ such that $b\neq c$ and $a= b\cdot c$.
			\end{enumerate}
		\end{enumerate}
		Moreover, property (b) implies that if  $a\in S_{n+1}\smallsetminus S_n$, and $b, a\cdot b\in S_n$, then  $\{b,a\cdot b\}$ is the unique pair given by (b) for  $a$.
	\end{prop}
	\begin{proof} 2 $\Rightarrow$ 1.  Let $N$ be a Steiner quasigroup and let $f:A\rightarrow N$ be a mapping. We extend $f$ to a homomorphism $\hat{f}: M\rightarrow N$.  We obtain $\hat{f}$ as the union $\hat{f}= \bigcup_{n<\omega}f_n$  of a chain  of partial homomorphisms $f_n : S_n \to N$. This means that  $f_n(a\cdot b)= f_n(a)\cdot f_n(b)$  whenever $a,b$ and $a\cdot b$  belong to $\mathrm{dom} (f_n)=S_n$. We start with $f_0=f$.  Recall that $A=S_0$ and that $a\cdot b\not\in S_0$  if $a,b\in S_0$ and, therefore, $f_0$ is a partial homomorphism with domain $S_0$. Now we extend  $f_n:S_n\rightarrow N$  to $f_{n+1}$. %We must define $f_{n+1}(a)$  for every $a\in S_{n+1}\smallsetminus S_n$. 
	For $a \in S_{n+1}\smallsetminus S_n$, there is a unique pair $\{b,c\}\subseteq S_n$ of distinct elements $b,c$ with $a=b\cdot c$. We let  $f_{n+1}(a) = f_n(b)\cdot f_n(c)$.  By (a) and (b), the only defined products involving $a$ in $S_{n+1}$ are   $a=b\cdot c=c\cdot b$, $a\cdot a=a$, $a\cdot b = b\cdot a= c$ and $a\cdot c= c\cdot a =b$ and $f_{n+1}$ respects these products.
		
		1 $\Rightarrow$ 2.  It is easy to construct a Steiner quasigroup $M^\prime\supseteq A $ with the properties described in 2. It is enough to start with $S^\prime_0=A$ and inductively define $S^\prime_{n+1}$ by adding to $S^\prime_n$ a different element $a\cdot b$ for each pair $\{a,b\}$ of elements of $S^\prime_n$ whose product is not defined in $S^\prime_n$. The only  products involving  the new element $c=a\cdot b$ in $S^\prime_{n+1}$ are $c=a\cdot b=b\cdot a$, $c\cdot c=c$, $c\cdot a= a\cdot c=b$ and $c\cdot b=b\cdot c= a$. We have proven that $A$ is a free base of $M^\prime= \bigcup_{n<\omega}S^\prime_n$. The uniqueness of free  Steiner quasigroups over a given set of generators guarantees that $M$ and $M^\prime$ are isomorphic over $A$.  The isomorphism can be used to obtain the levels $S_n$ in $M$ with the required properties.
	\end{proof}

	\begin{df}
		A \textbf{free ordering}  (or an \textbf{F-ordering})   of a Steiner quasigroup $M$ over $A\subseteq M$ is a linear ordering $<$ of $M$ such that for every $a\in M$  there is a unique pair $\{a_1,a_2\}\subseteq M$ with $a=a_1\cdot a_2$,  and for each $i$, either $a_i\in A$ or $a_i<a$. 
		
		Note that, when $a\in A$, the unique pair is given by $a_1=a_2=a$.
	\end{df}
	
	\begin{rmk}\label{May12_12.0} Let $M$ be a Steiner quasigroup with standard free construction $(S_n\mid n<\omega)$ over $A$. Let $<_n$ be a linear ordering of $S_n\smallsetminus S_{<n}$, and for $a,b\in M$ define
	$$a<b \Leftrightarrow \mbox{ for some } n, \ b\in S_n\smallsetminus S_{<n} \mbox{ and either } a\in S_{<n} \mbox{ or } a\in S_n\smallsetminus S_{<n} \mbox{ and } a<_n b.  $$ 
		Then $<$ is an F-ordering of $M$ over $A=S_0$, and $S_0$ is an initial segment. If every $<_n$ is a well-order, then $<$ is a well-order.
		\end{rmk}
		
	\begin{rmk}\label{May12_12.1}
		 Let $M$ be a Steiner quasigroup. If there is an F-ordering of $M$ over $A$, there is another one where $A$ is an initial segment. If the original order is a well-order, the new one is also a well-order and, moreover, $M=\langle A\rangle$.
	\end{rmk}	
	
	\begin{df} A \textbf{hyperfree ordering}  (or \textbf{HF-ordering}) of a partial Steiner triple system $A$ is a linear ordering $<$ of $A$ such that for every $a\in A$ there is at most one pair of smaller elements $\{a_1,a_2\}\subseteq A$ such that $a=a_1\cdot a_2$, that is, there is at most one block $\{a,a_1,a_2\}$ with $a_1,a_2<a$.
	\end{df}
The next proposition characterises free bases in terms of the existence of specific F-orderings and HF-orderings. Recall that a well-founded linear order is a well-order.
	\begin{prop}\label{May12_2.2} Assume $M$ is a Steiner quasigroup and $A\subseteq M$. The following are equivalent:
		\begin{enumerate}
			\item $A$ is a free base of $M$.
			\item There is an F-ordering of $M$ over $A$ which is a well-order.
			\item There is an HF-ordering of $M$ which is a well-order and $A$ is the set of points that are not a product of smaller elements.
		\end{enumerate}
		\end{prop}
		\begin{proof}
			1 $\Rightarrow$ 2. By Proposition~\ref{May12_2}, there is standard free construction $(S_n\mid n<\omega)$ of $M$ over $A$. Choose a well-order $<_n$ of every set $S_n\smallsetminus S_{<n}$. By Remark~\ref{May12_12.0}, there is an F-ordering $<$ of $M$ over $A$ which is a well-order extending all orders $<_n$.
			
			2 $\Rightarrow$ 3. By Remark~\ref{May12_12.1} there is an F-ordering $<$ of $M$ over $A$ which is a well-order and $A$ is an initial segment. We know that no element of $A$ is a product of smaller elements. If $a\in M\smallsetminus A$, then $a= a_1\cdot a_2$ where each $a_i$ is either smaller than $a$ or an element of $A$. Since $A$ is an initial segment, $a_1,a_2 <a$.
			
			3 $\Rightarrow$ 1.  Let $<$ be a well-order and an HF-ordering of $M$ and assume $A$ is the set of elements which are not a product of smaller elements. As in Remark~\ref{May12_12.1}, we can modify the order and make $A$ an initial segment. Hence, every element of $M\smallsetminus A$ is a product of smaller elements. By induction on the well-order it is easy to see that $M=\langle A\rangle$. Note that no element of $A$ is a product of two different elements of $A$. Let $N$ be a Steiner quasigroup and $f:A\rightarrow N$ a  mapping. By recursion on the well-order, it is easy to extend $f$ to some homomorphism $f^\prime:M\rightarrow N$.
		\end{proof}

\begin{prop}\label{2024_nov3} The following are equivalent for any partial Steiner triple system $A$:
\begin{enumerate}
\item $A$ is unconfined.
\item There is an HF-ordering of $A$.
\end{enumerate}	
\end{prop}	
\begin{proof} Consider first the case $A$ finite. By induction on $|A|$ it is easy to prove that, if $A$ is unconfined, then it has an HF-ordering. On the other hand, if there is an HF-ordering  of $A$, for any non-empty $A' \subseteq A$ the greatest element $a\in A'$ belongs to  at most one block of $A$, hence to at most a block of $A'$. Consider now the general case. If $A$ has an HF-ordering, every finite subset of $A$ has an HF-ordering and so it is unconfined, which implies that $A$ is unconfined. For the other direction, we use propositional compactness. Assume that every finite subset  of $A$ has an HF-ordering. For every pair  $(a,b)\in A\times A$, let $p_{ab}$ be a corresponding propositional variable. The set of sentences  
$$ \{p_{ab}\wedge p_{bc}\rightarrow p_{ac}\mid a,b,c\in A\}\cup \{\neg p_{aa}\mid a\in A\}\cup \{p_{ab}\vee p_{ba}\mid a,b\in A\}$$
 defines a linear ordering on $A$. If we add 
 $$\{\neg (p_{ba}\wedge p_{ca}\wedge p_{b^\prime a}\wedge p_{c^\prime a} )\mid \{a,b,c\}\neq\{a,b^\prime,c^\prime\} \text{ blocks of } A\}$$ 
 we are defining an HF-ordering on $A$. Every finite subset is satisfiable and compactness gives a valuation that satisfies the whole set.
\end{proof}

\begin{proof} Let $A$ be a finite set of generators of the unconfined Steiner quasigroup $M$. Let $(S_n\mid n<\omega)$ be the levels of $M$ over $A$. Every $S_n$ is finite. Each element of  $S_{n+1}\smallsetminus S_n$ is in at least one block intersecting $S_n$, and hence  there is an enumeration $(a_i\mid i<k)$ of $S_{n+1}\smallsetminus S_n$ such that  $\delta(S_na_1,\ldots,a_{i+1})\leq \delta(S_na_1,\ldots,a_i)\leq \delta(S_n)$. Concatenating  these enumerations after an arbitrary enumeration of $A=S_0$ gives an enumeration $(a_n\mid n<\omega)$ of $M$. Since $M$ is unconfined, $\delta(a_0,\ldots,a_n)\geq 0$  for every finite~$n$. Therefore there is some $n$  such that, for every $m\geq n$,\,  $\delta(a_0,\ldots,a_m)= \delta(a_0,\ldots,a_n)$. Since $M$ is unconfined, $A_m=\{a_i\mid i\leq m\}$ has an HF-ordering, and we may assume that the elements of $A_m$ which are not a product of smaller elements form an initial segment $A^\prime_m\subseteq A_m$. We start with this HF-ordering of $A_m$ and add the enumeration $(a_i\mid i>m)$, obtaining an F-ordering of $M$ over $A^\prime_m$ which is a well-order. Then $A^\prime_m$ is a free base of $M$ by Proposition~\ref{May12_2.2}.
	

\end{proof}				

\begin{thebibliography}{1}

\bibitem{andersenetal}
L.D. Andersen, A.J.W. Hilton, and E.~Mendelsohn.
\newblock Embedding partial {S}teiner triple systems.
\newblock {\em Proceedings of the London Mathematical Society}, 41:557--576,
  1980.

\bibitem{barcas3}
S.~Barbina and E.~Casanovas.
\newblock The elementary theory of free {S}teiner triple systems, 2024.
\newblock \url{http://arxiv.org/abs/2411.13723}.

\bibitem{fujiwara}
T.~Fujiwara.
\newblock Note on the isomorphism problem for free algebraic systems.
\newblock {\em Proc. Japan Acad.}, 32:662--664, 1956.

\bibitem{grishkovetal}
A.~Grishkov, D.~Rasskazova, M.~Rasskazova, and I.~Stuhl.
\newblock Free Steiner triple systems and their automorphism groups.
\newblock {\em Journal of Algebra and Its Applications}, 14(2) 1550025, 2015.

\bibitem{jontarski}
B.~J\'onsson and A.~Tarski.
\newblock On two properties of free algebras.
\newblock {\em Mathematica Scandinavica}, 9(1a):95--101, 1961.

\bibitem{lindner}
C.C. Lindner.
\newblock A survey of embedding theorems for {S}teiner systems.
\newblock {\em Annals of Discrete Mathematics}, 7:175--202, 1980.

\bibitem{sieben}
L.C. Siebenmann.
\newblock A characterization of free projective planes.
\newblock {\em Pacific Journal of Mathematics}, 15(1), 1965.

\bibitem{treash}
C.~Treash.
\newblock The completion of finite incomplete {S}teiner triple systems with
  applications to loop theory.
\newblock {\em Journal of Combinatorial Theory Series A}, 10, 1971.

\end{thebibliography}
\end{document}
